\section{Introduction}

Urban public transportation networks require continuous monitoring of passenger mobility to optimize fleet dispatching, evaluate route occupancy, and ensure passenger safety. While conventional automated passenger counting relies on optical cameras or infrared beams, these solutions incur substantial maintenance costs and raise passenger privacy concerns. In contrast, opportunistic Integrated Sensing and Communications (ISAC) transforms standard wireless infrastructure into dual-functional sensing nodes without dedicated radar hardware. Ubiquitous commodity Wi-Fi access points in transit vehicles continuously capture \ac{RSSI} streams from client mobile devices during routine network exchanges, enabling non-intrusive, privacy-preserving tracking at zero incremental hardware expense.

In public transport, mobility monitoring is cast as classifying four operational states: internal vehicle transit, alighting transitions, boarding events, and stationary platform waiting. Prior research established an initial benchmark collection of 2,251 \ac{RSSI} time-series samples \cite{matos2026icaisf}, on which a Gaussian Process classifier achieved an \ac{MCC} of 0.756 \cite{ribeiro2026eucnc}. However, supervised neural networks face severe limitations in dynamic vehicular environments. \ac{RSSI} streams suffer from non-stationary multipath fading and cross-traffic interference from concurrent transmitting devices. Supervised classifiers trained on limited annotations tend to memorize antenna-specific interference artifacts rather than underlying transit kinetics, resulting in brittle decision boundaries when operational conditions drift.

To overcome the fragility of supervised radio modeling, self-supervised learning via \ac{JEPA} world models provides a compelling framework. Rather than predicting raw noisy signals or reconstructing high-frequency radio fluctuations, \ac{JEPA} models predict abstract temporal representations in a latent embedding space, capturing macroscopic physical motion constraints while discarding non-semantic radio noise.

In this paper, we address these challenges through four primary contributions. First, we formulate passenger mobility sensing within a self-supervised world model framework, deploying four distinct \ac{JEPA} architectures (\ac{TS-JEPA}, LeJEPA, CF-JEPA, T-JEPA) alongside modern selective state-space models and regularized baselines. Second, we introduce a primary 3,511-sample transit sensing dataset, establishing that self-supervised world models generalize across collection regimes with \ac{MCC} improvements up to +0.042, whereas supervised classifiers suffer representational collapse. Third, we construct an 18-channel feature hierarchy, evaluate channel attribution via game-theoretic Shapley additive explanations (\ac{SHAP}), and demonstrate through empirical pruning that a compact 4-channel base subset preserves 99.4\% of peak classification performance (0.9143 vs. 0.9194 \ac{MCC}), enabling a 4.5-fold front-end computational reduction for real-time edge access points. Fourth, we develop temperature-calibrated stacked ensembles, achieving a peak \ac{MCC} of 0.9232 on the primary collection (0.9563 on multi-device noise) and isolating a compact three-member self-supervised ensemble delivering 0.9160 \ac{MCC} with only 3.42 million parameters.

The remainder of this paper is organized as follows. Section~II reviews related work in Wi-Fi sensing, ISAC architectures, and self-supervised sequence modeling. Section~III details the physical collection environment, the feature hierarchy, the \ac{JEPA} world models, and calibrated stacking. Section~IV analyzes empirical classification benchmarks, cross-collection transfer, feature importance, and edge computational trade-offs. Finally, Section~V concludes the paper.
